% Claude Marian-dogma reference: leaf-local typesetting primitives. This file
% owns only presentation. Every doctrinal and evidentiary judgment remains in
% the section files and research records.

\usepackage{endnotes}
% Long archival URLs in the references must break freely rather than
% overflowing the measure.
\usepackage{xurl}

\setcounter{tocdepth}{1}
\setlength{\emergencystretch}{2.2em}

\widowpenalty=10000
\clubpenalty=10000
\displaywidowpenalty=10000

% Work titles and transliterated technical terms are set through semantic
% commands so a later style change never requires editing the prose.
\newcommand{\work}[1]{\textit{#1}}
\newcommand{\term}[1]{\textit{#1}}

% The governing thesis is substantive content and may open the body.
\newcommand{\dogmathesis}[1]{%
  \begin{studybox}{Governing thesis}
  #1
  \end{studybox}}

% Defined-text block: the exact definitory formula with its authority, date,
% and locus in the frame title. The framed border marks quoted source text
% that is not project-owned expression.
\newenvironment{definedtext}[1]{%
  \begin{massbox}{Defined text --- #1}}{\end{massbox}}

% Endnotes carry citations and subordinate evidence detail; the heading joins
% the table of contents as an ordinary numbered section.
\renewcommand{\enoteheading}{\section{Notes}}
\renewcommand{\enoteformat}{%
  \noindent\footnotesize\makebox[1.8em][l]{\theenmark.}\hangindent=1.8em}

% Page-bounded synopsis card: one per dogma, kept together on a single page.
% sourcecard is unbreakable, so TeX moves an overfull card to a fresh page
% rather than splitting its fields.
\newenvironment{dogmasynopsis}[1]{%
  \begin{sourcecard}{Synopsis --- #1}\small}{\end{sourcecard}}
\newcommand{\synfield}[2]{\par\textbf{#1:} #2}

% Definitions table: the four dogmas set beside their defining or formulating
% act, the operative Latin wording as the identified witness prints it, and
% the level of the act with the assent it calls for.
\newenvironment{definitionstable}
{\footnotesize\begin{longtable}{>{\raggedright\arraybackslash}p{0.12\linewidth}>{\raggedright\arraybackslash}p{0.22\linewidth}>{\raggedright\arraybackslash}p{0.39\linewidth}>{\raggedright\arraybackslash}p{0.19\linewidth}}
\toprule
\textbf{Dogma} & \textbf{Act and date} & \textbf{Operative words (Latin witness)} & \textbf{Level and assent}\\
\midrule
\endhead
\bottomrule
\endfoot}
{\end{longtable}}

% Development-milestone table: one per dogma, carrying the date, the event,
% and the evidentiary register in which the event is claimed. The register
% column exists so a liturgical, legendary, or scholastic milestone can never
% be read as a definitional one.
\newenvironment{milestones}
{\footnotesize\begin{longtable}{>{\raggedright\arraybackslash}p{0.11\linewidth}>{\raggedright\arraybackslash}p{0.55\linewidth}>{\raggedright\arraybackslash}p{0.26\linewidth}}
\toprule
\textbf{Date} & \textbf{Milestone} & \textbf{Register}\\
\midrule
\endhead
\bottomrule
\endfoot}
{\end{longtable}}
